\begin{codebox}
\codeline{\textcolor{cmtgray}{\#\ ==============================}}
\codeline{\textcolor{cmtgray}{\#\ 4.\ NeOn\ 场景化方法（概览）}}
\codeline{\textcolor{cmtgray}{\#\ ==============================}}
\codeline{\textcolor{cmtgray}{\#\ NeOn（2009）不规定唯一流程，而是给出9种可组合的场景：}}
\codeline{\textcolor{cmtgray}{\#\ \ \ 场景1\ \ 从头构建（规格\(\rightarrow\)概念化\(\rightarrow\)实现）}}
\codeline{\textcolor{cmtgray}{\#\ \ \ 场景2\ \ 复用非本体资源（术语表/分类标准/数据库Schema\ \(\rightarrow\)\ 本体）}}
\codeline{\textcolor{cmtgray}{\#\ \ \ 场景3\ \ 复用本体资源（导入/裁剪已有本体）}}
\codeline{\textcolor{cmtgray}{\#\ \ \ 场景4\ \ 复用并重构（对已有本体做模块化改造）}}
\codeline{\textcolor{cmtgray}{\#\ \ \ 场景7\ \ 复用本体设计模式（ODP，如"部分-整体"模式）}}
\codeline{\textcolor{cmtgray}{\#\ \ \ ……}}
\codeline{\textcolor{cmtgray}{\#\ 工程现实中最常见组合：场景2\ +\ 场景3}}
\codeline{\textcolor{cmtgray}{\#\ （把企业已有的设备分类标准转成本体，同时对齐IOF顶层）}}
\codeblank{}
\end{codebox}
